\documentclass[11pt,a4paper,oneside]{report}

\usepackage[T1]{fontenc}
\usepackage[utf8]{inputenc}
\usepackage{lmodern}
\usepackage[a4paper,inner=3cm,outer=2.5cm,top=2.5cm,bottom=2.5cm]{geometry}
\usepackage{amsmath,amssymb,amsthm,mathtools}
\usepackage{bm}
\usepackage{booktabs}
\usepackage{enumitem}
\usepackage[numbers,sort&compress]{natbib}
\usepackage[hidelinks]{hyperref}
\usepackage[capitalise,noabbrev]{cleveref}

\setlength{\parskip}{0.3em}
\linespread{1.15}

% ---------------------------------------------------------------------
% Theorem environments (numbered within chapter, one shared counter)
% ---------------------------------------------------------------------
\theoremstyle{plain}
\newtheorem{theorem}{Theorem}[chapter]
\newtheorem{lemma}[theorem]{Lemma}
\newtheorem{proposition}[theorem]{Proposition}
\newtheorem{corollary}[theorem]{Corollary}
\theoremstyle{definition}
\newtheorem{definition}[theorem]{Definition}
\newtheorem{assumption}[theorem]{Assumption}
\newtheorem{example}[theorem]{Example}
\theoremstyle{remark}
\newtheorem{remark}[theorem]{Remark}

\crefname{assumption}{Assumption}{Assumptions}

% ---------------------------------------------------------------------
% Notation
% ---------------------------------------------------------------------
\newcommand{\R}{\mathbb{R}}
\newcommand{\E}{\mathbb{E}}
\newcommand{\X}{\mathcal{X}^\star}
\DeclarePairedDelimiter{\norm}{\lVert}{\rVert}
\DeclarePairedDelimiter{\abs}{\lvert}{\rvert}
\DeclarePairedDelimiterX{\inner}[2]{\langle}{\rangle}{#1, #2}
\DeclareMathOperator{\dist}{dist}
\DeclareMathOperator*{\argmin}{arg\,min}

\begin{document}

\setcounter{chapter}{2}

\chapter{Linear Convergence Without Convexity: the Polyak--{\L}ojasiewicz Condition}
\label{ch:pl}

\section{Introduction}
\label{sec:pl-intro}

The classical theory of first-order methods ties fast convergence to strong convexity. If $f\colon\R^d\to\R$ is $\mu$-strongly convex and has an $L$-Lipschitz gradient, gradient descent with step size $1/L$ reduces the optimality gap by a factor of $(1 - \mu/L)$ per iteration \citep{nesterov2004}. Many objectives that arise in practice are not strongly convex, however, and many are not convex at all: least squares with a rank-deficient design, logistic regression on separable data with a small ridge penalty, matrix factorisation, and the training losses of over-parameterised neural networks. Yet gradient methods are often observed to converge linearly on these problems.

This chapter studies one explanation for that behaviour. In 1963, Polyak \citep{polyak1963} showed that gradient descent converges linearly under a much weaker condition than strong convexity, namely that the squared gradient norm grows at least linearly with the suboptimality. The same inequality appears, as a special case of a more general one for real-analytic functions, in work of {\L}ojasiewicz \citep{lojasiewicz1963} published the same year. Following \citet{karimi2016} we call it the \emph{Polyak--{\L}ojasiewicz (PL) condition}.

Our aims are modest and are mainly to set up the tools used in Chapter~4. We (i)~state the PL condition and give examples of non-strongly-convex and non-convex functions that satisfy it (\cref{sec:pl-setting}); (ii)~show that it implies quadratic growth away from the solution set (\cref{sec:pl-qg}); (iii)~prove linear convergence of gradient descent (\cref{sec:pl-gd}); and (iv)~analyse stochastic gradient descent with constant and decreasing step sizes (\cref{sec:pl-sgd}). Most results in this chapter are known; our contribution is a self-contained presentation with complete proofs and explicit constants, which we rely on later.

\section{Setting and examples}
\label{sec:pl-setting}

Throughout, $\norm{\cdot}$ is the Euclidean norm. We consider the unconstrained problem
\begin{equation}\label{eq:problem}
  \min_{x\in\R^d} f(x),
\end{equation}
and write $f^\star = \inf_x f(x)$ and $\X = \{x : f(x) = f^\star\}$.

\begin{assumption}[Smoothness]\label{ass:smooth}
The function $f$ is differentiable, bounded below, and its gradient is $L$-Lipschitz:
$\norm{\nabla f(x) - \nabla f(y)} \le L\norm{x - y}$ for all $x, y\in\R^d$.
\end{assumption}

\begin{definition}[Polyak--{\L}ojasiewicz condition]\label{def:pl}
A differentiable function $f$ with $f^\star > -\infty$ satisfies the PL condition with constant $\mu > 0$ if
\begin{equation}\label{eq:pl}
  \frac{1}{2}\norm{\nabla f(x)}^2 \;\ge\; \mu\bigl(f(x) - f^\star\bigr) \qquad\text{for all } x\in\R^d.
\end{equation}
\end{definition}

The inequality says that every stationary point is a global minimiser: if $\nabla f(x) = 0$ then $f(x) = f^\star$. It does not say the minimiser is unique, and it does not require convexity.

\begin{lemma}\label{lem:mu-le-L}
If $f$ satisfies \cref{ass:smooth} and the PL condition with constant $\mu$, and $f$ is not constant, then $\mu \le L$.
\end{lemma}

\begin{proof}
By the descent lemma (\cref{lem:descent} below) applied with $y = x - \nabla f(x)/L$,
\[
  f^\star \le f\bigl(x - \tfrac1L\nabla f(x)\bigr) \le f(x) - \frac{1}{2L}\norm{\nabla f(x)}^2,
\]
so $\norm{\nabla f(x)}^2 \le 2L(f(x) - f^\star)$. Combining with \eqref{eq:pl} gives $2\mu(f(x) - f^\star) \le 2L(f(x) - f^\star)$, and choosing any $x$ with $f(x) > f^\star$ yields $\mu\le L$.
\end{proof}

We call $\kappa = L/\mu \ge 1$ the \emph{condition number} of $f$.

\begin{example}[Rank-deficient least squares]\label{ex:least-squares}
Let $A\in\R^{n\times d}$ be nonzero and $b\in\R^n$, and let $f(x) = \frac12\norm{Ax - b}^2$. If $A$ has a nontrivial null space then $f$ is not strongly convex, since it is constant along $\ker A$. Nevertheless, $f$ satisfies the PL condition with $\mu = \sigma_{\min}^{+}(A)^2$, the square of the smallest \emph{nonzero} singular value of $A$.
\end{example}

\begin{proof}[Verification]
Decompose $b = b_\parallel + b_\perp$ with $b_\parallel\in\operatorname{range}(A)$ and $b_\perp\in\ker A^\top$. Then $f^\star = \frac12\norm{b_\perp}^2$, and for any $x$,
\[
  f(x) - f^\star = \tfrac12\norm{Ax - b_\parallel}^2, \qquad
  \nabla f(x) = A^\top(Ax - b) = A^\top(Ax - b_\parallel),
\]
using $A^\top b_\perp = 0$. The vector $v = Ax - b_\parallel$ lies in $\operatorname{range}(A)$, which is the orthogonal complement of $\ker A^\top$, and on that subspace $\norm{A^\top v} \ge \sigma_{\min}^{+}(A)\norm{v}$. Hence
\[
  \tfrac12\norm{\nabla f(x)}^2 \ge \tfrac12\sigma_{\min}^{+}(A)^2\norm{v}^2 = \sigma_{\min}^{+}(A)^2\,\bigl(f(x) - f^\star\bigr). \qedhere
\]
\end{proof}

\begin{example}[A non-convex PL function]\label{ex:nonconvex}
The function $f(x) = x^2 + 3\sin^2 x$ on $\R$ is not convex, since $f''(x) = 2 + 6\cos 2x$ is negative near $x = \pi/2$. Its only stationary point is $x = 0$, and \citet{karimi2016} show that it satisfies \eqref{eq:pl} with $\mu = 1/32$.
\end{example}

\begin{example}[Composition with a linear map]\label{ex:composition}
If $g$ is $\mu$-strongly convex and $A$ is any matrix, then $f(x) = g(Ax)$ satisfies the PL condition, with a constant depending on $\mu$ and $\sigma_{\min}^+(A)$ \citep{karimi2016}. Least squares (\cref{ex:least-squares}) is the special case $g(z) = \frac12\norm{z - b}^2$. Logistic regression with a full-rank design restricted to a compact set is another.
\end{example}

\section{Quadratic growth}
\label{sec:pl-qg}

We first record the standard upper bound that comes with smoothness.

\begin{lemma}[Descent lemma]\label{lem:descent}
Under \cref{ass:smooth}, for all $x, y\in\R^d$,
\begin{equation}\label{eq:descent}
  f(y) \le f(x) + \inner{\nabla f(x)}{y - x} + \frac{L}{2}\norm{y - x}^2.
\end{equation}
\end{lemma}

\begin{proof}
By the fundamental theorem of calculus along the segment from $x$ to $y$,
\[
  f(y) - f(x) - \inner{\nabla f(x)}{y - x} = \int_0^1 \inner[\big]{\nabla f\bigl(x + t(y - x)\bigr) - \nabla f(x)}{y - x}\,dt.
\]
By Cauchy--Schwarz and the Lipschitz property, the integrand is at most $Lt\norm{y - x}^2$, and $\int_0^1 Lt\,dt = L/2$.
\end{proof}

The PL condition bounds the gradient from below in terms of function values. The next theorem converts it into a lower bound on function values in terms of the distance to the solution set. This \emph{quadratic growth} property is what makes the solution set well-conditioned: points with small suboptimality must be close to $\X$.

\begin{theorem}[PL implies quadratic growth]\label{thm:qg}
Suppose $f$ satisfies \cref{ass:smooth} and the PL condition with constant $\mu$. Then $\X$ is nonempty and closed, and
\begin{equation}\label{eq:qg}
  f(x) - f^\star \;\ge\; \frac{\mu}{2}\dist(x, \X)^2 \qquad\text{for all } x\in\R^d.
\end{equation}
\end{theorem}

\begin{proof}
Fix $x_0\in\R^d$. If $f(x_0) = f^\star$ there is nothing to prove, so assume $f(x_0) > f^\star$. Because $\nabla f$ is Lipschitz, the gradient flow
\[
  \dot{x}(t) = -\nabla f\bigl(x(t)\bigr), \qquad x(0) = x_0,
\]
has a unique solution defined for all $t\ge 0$. Every point of $\X$ is an equilibrium of this flow, so by uniqueness the trajectory never enters $\X$ in finite time, and $f(x(t)) > f^\star$ for all $t\ge0$.

\emph{Step 1: function values decay exponentially.} Along the flow,
\[
  \frac{d}{dt}\bigl(f(x(t)) - f^\star\bigr) = -\norm{\nabla f(x(t))}^2 \le -2\mu\bigl(f(x(t)) - f^\star\bigr),
\]
by \eqref{eq:pl}. Grönwall's inequality gives $f(x(t)) - f^\star \le e^{-2\mu t}(f(x_0) - f^\star)\to 0$.

\emph{Step 2: the path has finite length.} Let $g(t) = \sqrt{f(x(t)) - f^\star}$, which is differentiable since $f(x(t)) > f^\star$. Then
\[
  \dot{g}(t) = -\frac{\norm{\nabla f(x(t))}^2}{2\sqrt{f(x(t)) - f^\star}}
  \le -\sqrt{\frac{\mu}{2}}\,\norm{\nabla f(x(t))} = -\sqrt{\frac{\mu}{2}}\,\norm{\dot{x}(t)},
\]
where the inequality uses \eqref{eq:pl} in the form $\norm{\nabla f}/\sqrt{f - f^\star}\ge\sqrt{2\mu}$. Integrating from $0$ to $T$,
\begin{equation}\label{eq:path-length}
  \sqrt{\frac{\mu}{2}}\int_0^T\norm{\dot{x}(t)}\,dt \;\le\; g(0) - g(T) \;\le\; \sqrt{f(x_0) - f^\star}.
\end{equation}

\emph{Step 3: the flow converges to a point of $\X$.} By \eqref{eq:path-length} the total length of the trajectory is finite, so $x(t)$ is Cauchy as $t\to\infty$ and converges to some $\bar{x}$. By continuity and Step~1, $f(\bar{x}) = f^\star$, so $\bar{x}\in\X$; in particular $\X\neq\emptyset$. It is closed because $f$ is continuous.

\emph{Step 4: conclusion.} Since the straight-line distance is at most the path length, letting $T\to\infty$ in \eqref{eq:path-length} gives
\[
  \sqrt{\frac{\mu}{2}}\,\dist(x_0, \X) \le \sqrt{\frac{\mu}{2}}\,\norm{\bar{x} - x_0} \le \sqrt{f(x_0) - f^\star}.
\]
Squaring yields \eqref{eq:qg}.
\end{proof}

\begin{remark}\label{rem:hierarchy}
\citet{karimi2016} place the PL condition in a hierarchy of conditions used to prove linear convergence. For functions with Lipschitz gradient they show
\[
  \text{(SC)} \Rightarrow \text{(ESC)} \Rightarrow \text{(WSC)} \Rightarrow \text{(RSI)} \Rightarrow \text{(EB)} \equiv \text{(PL)} \Rightarrow \text{(QG)},
\]
where SC is strong convexity, ESC essential strong convexity, WSC weak strong convexity, RSI the restricted secant inequality, EB the Luo--Tseng error bound \citep{luo1993} and QG quadratic growth. \Cref{thm:qg} is the last implication. For convex functions, PL and QG are equivalent \citep{karimi2016,necoara2019}; without convexity QG is strictly weaker, since it allows non-optimal stationary points.
\end{remark}

\section{Gradient descent}
\label{sec:pl-gd}

\begin{theorem}[Linear convergence of gradient descent]\label{thm:gd}
Suppose $f$ satisfies \cref{ass:smooth} and the PL condition with constant $\mu$. Then the iterates of gradient descent with step size $1/L$,
\[
  x_{k+1} = x_k - \frac{1}{L}\nabla f(x_k),
\]
satisfy, for every $k\ge0$,
\begin{equation}\label{eq:gd-rate}
  f(x_k) - f^\star \le \Bigl(1 - \frac{\mu}{L}\Bigr)^{k}\bigl(f(x_0) - f^\star\bigr).
\end{equation}
\end{theorem}

\begin{proof}
Apply the descent lemma~\eqref{eq:descent} with $x = x_k$ and $y = x_{k+1} = x_k - \nabla f(x_k)/L$:
\[
  f(x_{k+1}) \le f(x_k) - \frac1L\norm{\nabla f(x_k)}^2 + \frac{1}{2L}\norm{\nabla f(x_k)}^2 = f(x_k) - \frac{1}{2L}\norm{\nabla f(x_k)}^2.
\]
By the PL inequality, $\norm{\nabla f(x_k)}^2 \ge 2\mu(f(x_k) - f^\star)$. Substituting and subtracting $f^\star$ from both sides,
\[
  f(x_{k+1}) - f^\star \le \Bigl(1 - \frac{\mu}{L}\Bigr)\bigl(f(x_k) - f^\star\bigr),
\]
and induction on $k$ gives \eqref{eq:gd-rate}. Note that $0\le 1 - \mu/L < 1$ by \cref{lem:mu-le-L}.
\end{proof}

\begin{corollary}[Iteration complexity]\label{cor:gd-complexity}
Under the conditions of \cref{thm:gd}, $f(x_k) - f^\star\le\varepsilon$ whenever
\[
  k \;\ge\; \kappa\,\log\frac{f(x_0) - f^\star}{\varepsilon}, \qquad \kappa = L/\mu.
\]
Moreover, $\dist(x_k, \X)^2 \le \frac{2}{\mu}(1 - 1/\kappa)^k(f(x_0) - f^\star)$.
\end{corollary}

\begin{proof}
The first claim follows from \eqref{eq:gd-rate} and $1 - 1/\kappa \le e^{-1/\kappa}$. The second follows by combining \eqref{eq:gd-rate} with \cref{thm:qg}.
\end{proof}

Two features of \cref{thm:gd} deserve comment. First, the rate matches the one for strongly convex functions exactly, with $\mu$ playing the role of the strong convexity constant; nothing is lost by dropping convexity. Second, the proof does not use the PL condition anywhere except along the iterates. It therefore suffices for \eqref{eq:pl} to hold on a set containing all the iterates, such as the sublevel set $\{x : f(x)\le f(x_0)\}$. This local version is what we verify for wide neural networks in Chapter~4, following the approach of \citet{liu2022}.

\section{Stochastic gradient descent}
\label{sec:pl-sgd}

In machine learning, $f$ is usually an average $f(x) = \frac1n\sum_{i=1}^n f_i(x)$ over training examples and exact gradients are too expensive. Stochastic gradient descent (SGD) replaces $\nabla f(x_k)$ by a cheap unbiased estimate $g_k$ \citep{robbins1951,bottou2018}:
\begin{equation}\label{eq:sgd}
  x_{k+1} = x_k - \alpha_k\,g_k.
\end{equation}
Let $\mathcal{F}_k$ denote the $\sigma$-algebra generated by $x_0, g_0, \dots, g_{k-1}$, so that $x_k$ is $\mathcal{F}_k$-measurable.

\begin{assumption}[Stochastic gradients]\label{ass:noise}
For every $k$, $\E[g_k\mid\mathcal{F}_k] = \nabla f(x_k)$ and $\E\bigl[\norm{g_k - \nabla f(x_k)}^2\bigm|\mathcal{F}_k\bigr] \le \sigma^2$.
\end{assumption}

Write $\delta_k = \E[f(x_k)] - f^\star$. The analysis rests on a one-step recursion.

\begin{lemma}[One-step progress]\label{lem:sgd-step}
Under \cref{ass:smooth,ass:noise} and the PL condition, if $0 < \alpha_k \le 1/L$ then
\begin{equation}\label{eq:sgd-recursion}
  \delta_{k+1} \le (1 - \alpha_k\mu)\,\delta_k + \frac{L\alpha_k^2\sigma^2}{2}.
\end{equation}
\end{lemma}

\begin{proof}
The descent lemma with $y = x_{k+1}$ gives
\[
  f(x_{k+1}) \le f(x_k) - \alpha_k\inner{\nabla f(x_k)}{g_k} + \frac{L\alpha_k^2}{2}\norm{g_k}^2.
\]
Taking conditional expectations and using \cref{ass:noise}, together with the bias--variance identity $\E[\norm{g_k}^2\mid\mathcal{F}_k] = \norm{\nabla f(x_k)}^2 + \E[\norm{g_k - \nabla f(x_k)}^2\mid\mathcal{F}_k]$,
\[
  \E[f(x_{k+1})\mid\mathcal{F}_k] \le f(x_k) - \alpha_k\Bigl(1 - \frac{L\alpha_k}{2}\Bigr)\norm{\nabla f(x_k)}^2 + \frac{L\alpha_k^2\sigma^2}{2}.
\]
Since $\alpha_k\le1/L$ we have $1 - L\alpha_k/2 \ge 1/2$, and by the PL inequality $\frac12\norm{\nabla f(x_k)}^2\ge\mu(f(x_k) - f^\star)$. Therefore
\[
  \E[f(x_{k+1})\mid\mathcal{F}_k] - f^\star \le (1 - \alpha_k\mu)\bigl(f(x_k) - f^\star\bigr) + \frac{L\alpha_k^2\sigma^2}{2}.
\]
Taking total expectations completes the proof.
\end{proof}

\begin{theorem}[SGD with constant step size]\label{thm:sgd-constant}
Under the conditions of \cref{lem:sgd-step} with $\alpha_k = \alpha\in(0, 1/L]$ for all $k$,
\begin{equation}\label{eq:sgd-constant}
  \delta_k \le (1 - \alpha\mu)^k\,\delta_0 + \frac{L\alpha\sigma^2}{2\mu}.
\end{equation}
\end{theorem}

\begin{proof}
Unrolling \eqref{eq:sgd-recursion},
\[
  \delta_k \le (1 - \alpha\mu)^k\delta_0 + \frac{L\alpha^2\sigma^2}{2}\sum_{j=0}^{k-1}(1 - \alpha\mu)^j
  \le (1 - \alpha\mu)^k\delta_0 + \frac{L\alpha^2\sigma^2}{2}\cdot\frac{1}{\alpha\mu}. \qedhere
\]
\end{proof}

With a constant step size, SGD converges linearly, but only to a neighbourhood of the optimum whose size is proportional to $\alpha\sigma^2$. Halving the step size halves the neighbourhood but also halves the rate. When $\sigma = 0$ at the optimum, as in over-parameterised models that interpolate the data, the neighbourhood vanishes and SGD converges linearly to the optimum \citep{vaswani2019,gower2019}; this is the regime of Chapter~4. In general, reaching arbitrary accuracy requires a decreasing step size.

\begin{theorem}[SGD with decreasing step size]\label{thm:sgd-decreasing}
Under the conditions of \cref{lem:sgd-step}, let $\alpha_k = \beta/(\gamma + k)$ with $\beta > 1/\mu$ and $\gamma > 0$ chosen so that $\alpha_0 = \beta/\gamma \le 1/L$. Then for all $k\ge0$,
\begin{equation}\label{eq:sgd-decreasing}
  \delta_k \le \frac{\nu}{\gamma + k}, \qquad\text{where}\quad
  \nu = \max\Bigl\{\frac{\beta^2 L\sigma^2}{2(\beta\mu - 1)},\; \gamma\,\delta_0\Bigr\}.
\end{equation}
\end{theorem}

\begin{proof}
We argue by induction; the case $k = 0$ holds by the definition of $\nu$. Suppose \eqref{eq:sgd-decreasing} holds for some $k$ and write $\hat{k} = \gamma + k$. The step sizes are decreasing, so $\alpha_k\le\alpha_0\le1/L$ and \cref{lem:sgd-step} applies:
\begin{align*}
  \delta_{k+1}
  &\le \Bigl(1 - \frac{\beta\mu}{\hat{k}}\Bigr)\frac{\nu}{\hat{k}} + \frac{\beta^2 L\sigma^2}{2\hat{k}^2}
   = \frac{\hat{k} - 1}{\hat{k}^2}\,\nu - \frac{(\beta\mu - 1)\nu - \beta^2 L\sigma^2/2}{\hat{k}^2} \\
  &\le \frac{\hat{k} - 1}{\hat{k}^2}\,\nu
   \;\le\; \frac{\nu}{\hat{k} + 1},
\end{align*}
where the second inequality uses $\nu\ge\beta^2L\sigma^2/(2(\beta\mu - 1))$, and the last uses $(\hat{k} - 1)(\hat{k} + 1)\le\hat{k}^2$. Since $\hat{k} + 1 = \gamma + (k+1)$, this completes the induction. (In the first line we used $1 - \beta\mu/\hat{k}\ge 0$, which holds because $\beta\mu/\hat{k} = \alpha_k\mu\le\mu/L\le1$.)
\end{proof}

\Cref{thm:sgd-decreasing} gives the $O(1/k)$ rate that is optimal for stochastic first-order methods on strongly convex problems \citep{moulines2011,bottou2018}, now under the PL condition alone. The requirement $\beta > 1/\mu$ is not cosmetic: if the step sizes decay too quickly relative to the curvature, the rate degrades to $O(k^{-\beta\mu})$.

\begin{table}[t]
  \centering
  \caption{Summary of the convergence guarantees proved in this chapter. All assume \cref{ass:smooth} and the PL condition with constant $\mu$; $\kappa = L/\mu$.}
  \label{tab:pl-summary}
  \begin{tabular}{@{}llll@{}}
    \toprule
    Method & Step size & Guarantee & Result \\
    \midrule
    Gradient flow & --- & $f - f^\star\le e^{-2\mu t}\,(f(x_0) - f^\star)$ & \cref{thm:qg}, Step 1 \\
    Gradient descent & $1/L$ & $(1 - 1/\kappa)^k$ contraction & \cref{thm:gd} \\
    SGD & constant $\alpha\le 1/L$ & linear to $O(L\alpha\sigma^2/\mu)$ ball & \cref{thm:sgd-constant} \\
    SGD & $\beta/(\gamma + k)$ & $O(1/k)$ & \cref{thm:sgd-decreasing} \\
    \bottomrule
  \end{tabular}
\end{table}

\section{Discussion}
\label{sec:pl-discussion}

\Cref{tab:pl-summary} summarises the results. The common thread is that every proof uses the PL inequality only to convert a decrease in $\norm{\nabla f}^2$, which the descent lemma provides, into a decrease in $f - f^\star$, which is what we want. Convexity never enters.

Two limitations matter for what follows. First, the PL condition is global, and it is rarely possible to verify it globally for non-convex models. The Kurdyka--{\L}ojasiewicz framework \citep{kurdyka1998,bolte2007,attouch2010} relaxes the exponent and localises the inequality, at the cost of rates that depend on unknown constants. Second, the constant $\mu$ is usually unknown, so step-size schedules such as the one in \cref{thm:sgd-decreasing} cannot be implemented exactly. In Chapter~4 we address the first issue for wide two-layer networks by showing that a local PL condition holds, with high probability over the initialisation, on a ball that contains the entire optimisation trajectory. The second issue we leave to adaptive methods, discussed in Chapter~6.

\bibliographystyle{plainnat}
\bibliography{06-references}

\end{document}
